\begin{abstract}
Event-related potentials (ERPs) provide high-temporal-resolution measurements of neural response dynamics, but their use in biomedical data analysis is limited by subject variability, low signal-to-noise ratio, and sensitivity to acquisition and preprocessing perturbations. This paper evaluates a fixed spiking reservoir as a nonlinear temporal encoding operator for affective ERP observations. Multichannel ERP traces are transformed by a leaky integrate-and-fire reservoir, summarized using six post-onset spike-count bins, and compressed into a subject-level representation for affective condition analysis. The study tests whether this binned spiking code provides a stable temporal representation under controlled perturbations, rather than claiming classifier superiority over conventional ERP features. Evaluation uses subject-grouped validation, balanced accuracy, macro-F1, classical ERP baselines, and perturbation tests involving additive amplitude noise, temporal jitter, and channel dropout. Results indicate that the reservoir embedding is not a uniformly stronger static classifier than classical ERP-derived features, but it exhibits a distinct perturbation-response profile. The findings support a bounded interpretation of spiking reservoirs as perturbation-characterized temporal measurement operators for biomedical ERP signal analysis.
\end{abstract}

\begin{IEEEkeywords}
EEG, event-related potentials, spiking neural networks, reservoir computing, biomedical signal processing, affective computing, perturbation analysis.
\end{IEEEkeywords}
